\documentclass[11pt]{article}
\usepackage[a4paper,margin=24mm]{geometry}
\usepackage[T1]{fontenc}
\usepackage{amsmath,amssymb,amsthm,mathtools,xurl}
\usepackage[hidelinks]{hyperref}
\allowdisplaybreaks
\setlength{\emergencystretch}{3em}
\newcommand{\R}{\mathbb R}
\newcommand{\C}{\mathbb C}
\newcommand{\norm}[1]{\left\lVert#1\right\rVert}
\newcommand{\lp}[1]{\log^+\!#1}
\title{Moving-interval native margins and finite raw-moment precision}
\author{}
\date{20 September 2026}
\begin{document}\maketitle

This independently derives the native precision part NP1--22 of the
incoming continuation. The complete interval inverse, all constants,
the original-coordinate matrix floor, the finite convolution map and
both types of receiving error are given below. It strengthens the
earlier native precision proof \cite{oldNP}, the complex-product
composition \cite{MC}, and the explicit resolvent degree bound
\cite{NR}. The original full polynomial source and its attained
quotients are retained. The estimates do not evaluate the still
unevaluated arithmetic moment centres or their signed currents.

\section{Objects, domains and established original envelopes}

Fix the existing full quartet domain
\begin{equation}\tag{MI1}\label{mi:objects}
\begin{gathered}
 \rho_*=\tfrac12+\delta+i\gamma,\quad0<\delta<\tfrac12,\quad
 \gamma>2,\quad m\in\mathbb Z_{\ge1},\quad g=2\xi,\\
 h(z)=\prod_{\epsilon,\eta=\pm1}
       (z-\tfrac12-\epsilon\delta-i\eta\gamma)^m,\qquad
 w_h(y)=\frac{|(g/h)(1/2+iy)|^2}{2\pi},\\
 m_k=w_h^{*k},\quad k=4\ell+1\ge9,\quad
 e_k=1+k(m-1),\quad q=e_k(k+1)^2,\quad c=k/2,\\
 \chi_k(S)=\prod_{a,b=0}^k
 [S-c-(2a-k)\delta-i(2b-k)\gamma]^{e_k}.
\end{gathered}
\end{equation}
This retains the stipulated actual quartet and its full multiplicity;
it asserts no existence of an off-line zero. The physical coordinate
is still $S=c+iy$. The original Gamma comparison, with its actual
mass $\int\sigma(y)\,dy=\sqrt{2\pi}$, is
\begin{equation}\tag{MI2}\label{mi:env}
\begin{gathered}
 \sigma(y)=|\Gamma(1/4+iy/2)|^2/(2\pi),\quad
 a_k(1+y^2)^{-M}\sigma(y)\le m_k(y)\le A_k\sigma(y),\\
 c_\Gamma e^{-\alpha|y|}(1+|y|)^{-1/2}
 \le\sigma(y)\le C_\Gamma e^{-\alpha|y|}(1+|y|)^{-1/2},\\
 \alpha=\pi/2,\quad M=21+4m,\quad p=8m-9/2,\quad B_h=42+8m,\\
 c_\Gamma=\sqrt2e^{-7/6},\quad C_\Gamma=\sqrt{2\sqrt5}e^{2/3},\\
 \vartheta_h=\int_{-1}^1w_h(y)\,dy,\quad
 M_\alpha=\int_\R e^{\alpha y}w_h(y)\,dy,\\
 a_k=\frac{c_h\vartheta_h^{k-3}e^{-\alpha(k-3)}(k-2)^{-B_h}}
                 {C_\Gamma2^{B_h/2}},\qquad
 A_k=\frac{C_h^{\rm tilt}}{c_\Gamma}k^{p+1}M_\alpha^{k-1}.
\end{gathered}
\end{equation}
These are the established CAI16--17, TW7 and BT12--23 envelope
constants as used in \cite{oldNP,BT}. They are never assigned
artificial arithmetic values. The exact parity measures are
\begin{equation}\tag{MI3}\label{mi:parity}
 d\nu_b(x)=x^{b-1/2}m_k(\sqrt x)\,dx,\quad
 d\widetilde\nu_b=s\,d\nu_b,\quad
 s(x)=\frac{\tanh(\alpha\sqrt x)}{\sqrt x},\quad b=0,1.
\end{equation}
For either measure $\lambda$, put
$H_{\lambda,t}=[\int x^{i+j}\,d\lambda]_{0\le i,j\le t}$.
These are full original integrals, in increasing monomials.

\section{An exact interval inverse in those monomials}

For $R>1$, let
$J_{t,R}[i,j]=(R^{i+j+1}-1)/(i+j+1)$.
Use the classical Legendre polynomials $P_n$; the Rodrigues formula
is the human result recorded in Dunster's DLMF 14.7.13
\cite{Dunster}. In the variable $u$ it gives
\[
 P_n(2u-1)=\frac1{n!}\frac{d^n}{du^n}[u^n(u-1)^n]
 =\sum_{l=0}^n(-1)^{n-l}\binom nl\binom{n+l}l u^l.
\]
This equality follows by expanding $(u-1)^n$ before differentiating.
Integration by parts $n$ times makes this polynomial orthogonal on
$[0,1]$ to every polynomial of degree less than $n$: all boundary
terms vanish because $u^n(u-1)^n$ has zeros of order $n$ at both
endpoints. Its squared norm is
\[
 \frac{(2n)!}{(n!)^2}\int_0^1u^n(1-u)^n\,du=\frac1{2n+1}.
\]
For the last integral, integration by parts repeatedly gives
$\int_0^1u^n(1-u)^n du=(n!)^2/(2n+1)!$.
Define the actual coefficient matrix
\begin{equation}\tag{MI4}\label{mi:C}
\begin{split}
 C_{jn}(R)
 &=[x^j]P_n\!\left(\frac{2x-R-1}{R-1}\right)\\
 &=(-1)^{n-j}\sum_{l=j}^n
   \binom nl\binom{n+l}l\binom lj(R-1)^{-l}
       \qquad(0\le j\le n\le t).
\end{split}
\end{equation}
Every diagonal entry is nonzero. The substitution
$u=(x-1)/(R-1)$ has Jacobian $R-1$; hence, with both coordinate
maps retained,
\begin{equation}\tag{MI5}\label{mi:inverse}
 \begin{gathered}
 C^TJ_{t,R}C=\operatorname{diag}_{n=0}^t\frac{R-1}{2n+1},\\
 J_{t,R}^{-1}
 =C\operatorname{diag}_{n=0}^t\frac{2n+1}{R-1}C^T,\qquad
 J_{t,R}\succeq S_t(R)^{-1}I,\\
 S_t(R)=\frac1{R-1}\sum_{n=0}^t(2n+1)\sum_{j=0}^n C_{jn}(R)^2.
 \end{gathered}
\end{equation}
The last assertion follows from
$\lambda_{\max}(J^{-1})\le\operatorname{tr}(J^{-1})=S_t(R)$.
In particular at $R=2$, the inverse traces for $t=0,\ldots,6$ are
\[
 1,\ 40,\ 2685,\ 211936,\ 17845105,\ 1556000328,\ 138680834005.
\]
The coefficient signs in \eqref{mi:C} alternate. Evaluation at
$x=-1$, together with $P_n(-z)=(-1)^nP_n(z)$, therefore gives
$\sum_j|C_{jn}(R)|=P_n(z_R)$, where
$z_R=(R+3)/(R-1)>1$.

For $z>1$ the Legendre generating formula of
\cite{DLMF18} gives the integral identity
\[
 P_n(z)=\frac1\pi\int_0^\pi
        (z+\sqrt{z^2-1}\cos\theta)^n\,d\theta.
\]
To verify the identity directly from that generating formula, sum
the geometric series for sufficiently small real $v$ inside the
integral. The substitution $u=\tan(\theta/2)$ evaluates the resulting
integral as
$[(1-vz)^2-v^2(z^2-1)]^{-1/2}=(1-2zv+v^2)^{-1/2}$.
Coefficient comparison proves the assertion. Thus, setting
$\Lambda_R=z_R+\sqrt{z_R^2-1}$, one gets
$P_n(z_R)\le\Lambda_R^n$ and
\begin{equation}\tag{MI6}\label{mi:intervalclosed}
 S_t(R)\le\frac{(t+1)^2}{R-1}\Lambda_R^{2t},
 \qquad J_{t,R}\succeq
 \frac{R-1}{(t+1)^2}\Lambda_R^{-2t}I.
\end{equation}
The sum $\sum_{n=0}^t(2n+1)=(t+1)^2$ accounts for every factor.

\section{A moving interval in the full original native integral}

For $x\in[1,R]$, \eqref{mi:env} gives a common density floor for
both $\nu_b$:
\begin{equation}\tag{MI7}\label{mi:d}
 d_k(R)=\frac{a_kc_\Gamma e^{-\alpha\sqrt R}}
 {(1+R)^M\sqrt R\sqrt{1+\sqrt R}}.
\end{equation}
The $b=1$ density has the extra factor $x\ge1$. The function $s$
decreases: for $u>0$ the derivative of
$\tanh u-u\operatorname{sech}^2u$ is
$2u\operatorname{sech}^2u\tanh u>0$, and its value at zero is zero.
For every coefficient vector, restrict only its nonnegative squared
integral to $[1,R]$. Equations \eqref{mi:inverse}--\eqref{mi:d} give
\begin{equation}\tag{MI8}\label{mi:floor}
 \begin{gathered}
 H_{\nu_b,t}\succeq h_{k,t}(R)I,\qquad
 H_{\widetilde\nu_b,t}\succeq s(R)h_{k,t}(R)I,\\
 h_{k,t}(R)=d_k(R)\frac{R-1}{(t+1)^2}\Lambda_R^{-2t}.
 \end{gathered}
\end{equation}
One may replace $h_{k,t}(R)$ by the larger exact expression
$d_k(R)/S_t(R)$. No source matrix has become an interval matrix.

Take $L=t+1\ge1$ and $R_t=1+4L$.
Since $\cosh u\ge1+u^2/2$, one has
$\operatorname{arcosh}(1+v)\le\sqrt{2v}$. It follows that
\[
 2t\log\Lambda_{R_t}\le2\sqrt2\sqrt L,\qquad
 \alpha\sqrt{R_t}\le\alpha+2\alpha\sqrt L.
\]
Also $1+R_t\le6L$, $\sqrt{R_t}\le\sqrt5\sqrt L$ and
$1+\sqrt{R_t}\le(1+\sqrt5)\sqrt L$. Substitution in
\eqref{mi:floor}, including $(R_t-1)/L^2=4/L$, proves
\begin{equation}\tag{MI9}\label{mi:rootfloor}
 \begin{gathered}
 H_{\nu_b,t}\succeq f_{k,t}I,\qquad
 H_{\widetilde\nu_b,t}\succeq\widetilde f_{k,t}I,\\
 f_{k,t}=
 \frac{4a_kc_\Gamma e^{-\alpha}}
 {6^M\sqrt{5(1+\sqrt5)}}L^{-(M+7/4)}
 e^{-(\pi+2\sqrt2)\sqrt L},\\
 \widetilde f_{k,t}=\frac{\tanh\alpha}{\sqrt5\sqrt L}f_{k,t}.
 \end{gathered}
\end{equation}
The tilted factor follows from
$s(R_t)\ge\tanh\alpha/\sqrt{R_t}$, since $R_t\ge1$.
All these statements include $t=0$.

Here is an explicit logarithmic form of this estimate. Put
\[
 D_h=\frac{4c_\Gamma e^{-\alpha}}
 {6^M\sqrt{5(1+\sqrt5)}},\quad
 E_h=\log^+\frac{C_\Gamma2^{B_h/2}}{c_h},\quad
 v_h=\log^+(\vartheta_h^{-1}),\quad
 \log^+u=\max(0,\log u).
\]
Then
\begin{equation}\tag{MI10}\label{mi:log}
\begin{split}
 \log^+(f_{k,t}^{-1})\le{}&
 \log^+(D_h^{-1})+E_h+(k-3)(\alpha+v_h)
       +B_h\log(k-2)\\
 &+(M+7/4)\log L+(\pi+2\sqrt2)\sqrt L.
\end{split}
\end{equation}
For $\widetilde f$ add
$\log(\sqrt5/\tanh\alpha)+\frac12\log L$.
If $0\le t\le Cq$ for a fixed $C>0$, then
$L\le(C+1)q$ and $k\le\sqrt q$. Consequently
\[
 \log^+(f_{k,t}^{-1}),\
 \log^+(\widetilde f_{k,t}^{-1})
 =O_{h,C}(\sqrt q).
\]
This is the cost of the explicit lower certificate. It assigns no
asymptotic equality to a native eigenvalue.

\section{Weighted inverses, the Radau denominator and the new degree bound}

Let $w(x)=\sum_lw_lx^l$ with $w_l\ge0$ and $w(1)>0$, and let
$W_t(w)=[\sum_lw_l\mu_{i+j+l}]_{i,j=0}^t$ for an actual measure
$\lambda$. Choose any proved floor $h_{\lambda,t}$ from
\eqref{mi:floor}, its inverse-trace improvement, \eqref{mi:rootfloor},
or the old fixed-interval bound; their maximum is again a floor.
On the selected interval $w(x)\ge w(1)$, so
\begin{equation}\tag{MI11}\label{mi:weighted}
 W_t(w)\succeq w(1)h_{\lambda,t}I,\qquad
 \norm{W_t(w)-\widehat W_t(w)}
 \le(t+1)w(1)\varepsilon.
\end{equation}
The second statement follows by summing the absolute entry errors
in each row when every required scalar moment error is at most
$\varepsilon$. Therefore
$\varepsilon\le h_{\lambda,t}/[4(t+1)]$ gives a positive lower
outward endpoint. The same inverse and quadratic-expression estimates
in \cite{oldNP} follow by substituting this proved larger floor.
Any additional $\varepsilon\le1$ guard in those estimates is retained.
This covers $x^l,x+a,x(x+a),x(x+a)(x+b)$ with their actual positive
poles and shifts; all coefficient sums are exactly their $w(1)$.

For the actual $r$-dimensional tilted Radau system, let
$H=[\mu_{i+j}]_{i,j<r}$, $D=[\mu_{i+j+1}]_{i,j<r}$,
$e=(0,\ldots,0,1)^T$, $\kappa=(e^*D^{-1}e)^{-1}$ and
$W_G=D+aH$. The attained coefficient minimum gives
$D-\kappa ee^*\succeq0$. Hence
$W_G-\kappa ee^*\succeq a h_{\widetilde\nu,r-1}I$.
Congruence by $W_G^{-1/2}$ proves
\begin{equation}\tag{MI12}\label{mi:radau}
 1-\kappa e^*W_G^{-1}e
 \ge\frac{a h_{\widetilde\nu,r-1}}{U_D+aU_H}>0,
\end{equation}
where $U_D,U_H$ are the original positive trace caps.
Indeed the only eigenvalue of $I-\kappa uu^*$ which can differ
from one is $1-\kappa u^*u$, and
$\lambda_{\max}(W_G)\le U_D+aU_H$.

This also feeds the same native resolvent degree formula NR19.
At original total degree $N$, put $n_b=\lfloor(N-b)/2\rfloor$,
omitting empty blocks, and
$H_{\rm base}=(2/\pi)\operatorname{diag}(H_{\widetilde\nu_b,n_b})$.
For $a_j=4j^2$ and the original padded evaluation vector $z_{b,j}$,
\begin{equation}\tag{MI13}\label{mi:newchi}
 \chi_{b,j}=z_{b,j}^*H_{\rm base}^{-1}z_{b,j}
 \le \frac{\pi}{2h_{\widetilde\nu_b,n_b}}
          \sum_{l=0}^{n_b}a_j^{2l}.
\end{equation}
It follows by inverting the block floor; no evaluation direction is
changed. The minimum of this bound and the earlier NR26 Gamma
kernel bound is another proved upper bound on the same scalar.
Using it in
\[
 r_{b,j}=\max\left\{n_b+1,\
 \left\lceil\left(\frac{4a_j\mathcal K_{k,j}\chi_{b,j}^{\rm upper}}
 {\pi\varepsilon_{b,j}}\right)^{1/(2j)}\right\rceil-1\right\},\qquad
 \mathcal K_{k,j}=\frac{\pi\sqrt{2\pi}}{4j}A_k((5/4)_{j-1})^2,
\]
preserves the already proved same-direction inner matrix guarantee.

\section{Exact exponential-generating-function transport}

Define actual onefold and convolution moments by
$\eta_j=\int y^jw_h(y)\,dy$ and
$\zeta_j=\int y^jm_k(y)\,dy$.
For every integer $L\ge0$ the finite polynomial
$A_L(z)=\sum_{j=0}^L\eta_jz^j/j!$ satisfies
\begin{equation}\tag{MI14}\label{mi:egf}
 \zeta_j=j![z^j]A_L(z)^k,\qquad 0\le j\le L.
\end{equation}
Indeed insert the finite multinomial expansion of
$(y_1+\cdots+y_k)^j$ in the original convolution integral.
Fubini is valid because all absolute moments are finite by the
original exponential bound. Coefficients of degree at most $L$
cannot involve a higher input degree. In particular
$\zeta_0=\eta_0^k$ with the complete original mass.

The original BT14--16 estimate \cite{BT} gives
$e^{\alpha|y|}w_h(y)\le C_h^{\rm tilt}(1+|y|)^{-p}$.
Thus both of the following are explicit fixed-packet caps:
\begin{equation}\tag{MI15}\label{mi:B}
 \begin{gathered}
 M_h(1)=\int_\R\cosh(y)w_h(y)\,dy
 \le M_\alpha\le B_*:=\frac{2C_h^{\rm tilt}}{p-1},\\
 M_h(1)\le B_{\rm in}:=\frac{2C_h^{\rm tilt}}{\alpha-1},
 \qquad B_*<B_{\rm in}.
 \end{gathered}
\end{equation}
The first line uses evenness, monotonicity of $\cosh$ for positive
arguments, and
$\int_\R(1+|y|)^{-p}dy=2/(p-1)$.
The second follows alternatively by integrating
$C_h^{\rm tilt}e^{-(\alpha-1)|y|}$.
Here $p\ge7/2>\alpha>1$, proving the strict comparison.
The incoming $B_{\rm in}$ is therefore valid; $B_*$ improves it.
The unchanged fully explicit BT12 constant is
\[
\begin{gathered}
 R_0=\sqrt{\delta^2+\gamma^2},\quad T_h=\max(1,2R_0),\
 Z_*=2+\sqrt{5/2},\quad\beta=p-1/2,\\
 C_h^{\rm tilt}=\max\left\{
 \frac{C_\Gamma}{\sqrt\pi}(T_h^2+1/4)^2
 [1+2\sqrt{T_h^2+1/4}]^2\delta^{-8m}(1+T_h)^\beta,\
 \frac{25C_\Gamma Z_*^2}{16\sqrt\pi}
 (4/3)^{4m}(1+1/T_h)^\beta\right\}.
\end{gathered}
\]
Its compact and tail proof is BT3--14 of \cite{BT}; this
calculation uses that existing proved original bound.

Odd moments vanish, and even moments are nonnegative. Thus
$\sup_{|z|\le1}|A_L(z)|\le M_h(1)\le B$ for either cap
$B=B_*$ or $B=B_{\rm in}$. If
$|\widehat\eta_j-\eta_j|\le\tau\le1/3$ for all $j\le L$, then
\[
 \sup_{|z|\le1}|\widehat A_L-A_L|
 \le\tau\sum_{j=0}^L1/j!<3\tau\le1.
\]
Telescope the difference of the two $k$th powers to bound it by
$3k\tau(B+1)^{k-1}$ throughout the unit circle.
Integrating that polynomial against $e^{-ij\theta}$ on this
circle isolates its $j$th coefficient and proves
\begin{equation}\tag{MI16}\label{mi:egferror}
 |\widehat\zeta_j-\zeta_j|
 \le3k\tau(B+1)^{k-1}j!,\qquad 0\le j\le L.
\end{equation}
Consequently all untilted source entries through total degree $N$
are simultaneously certified by the single tolerance
\begin{equation}\tag{MI17}\label{mi:tau}
 \tau\le\min\left\{\frac13,\
 \frac{\varepsilon}{3k(B+1)^{k-1}(2N)!}\right\}.
\end{equation}
Both parity blocks use the same $\widehat\zeta_{2r}$ sequence and
its literal index shift. The assertion concerns the untilted source
moments. Multiplying by the nonpolynomial $s(y^2)$ is not represented
by a finite coefficient in \eqref{mi:egf}; the earlier resolvent
receiver retains that separate exact operation.

\section{The same original quotient and complex mixed product}

Here use exactly the proved simple-quartet current domain
$m=1$, $k\equiv1\pmod4$, $k\ge77$, $q=(k+1)^2$, $N\ge q-1$.
Let $H$ be the direct sum of the actual untilted parity Grams,
$d=\max_b(n_b+1)$, and let $h_{\min}$ be the minimum of their
improved floors above. Retain the original source-to-packet map
\begin{equation}\tag{MI18}\label{mi:map}
 \begin{gathered}
 J=J_{S,N}L_N^{-1},\quad L_N=\Pi_NT_N,\quad
 (T_N)_{rj}=\binom jr c^{j-r}i^r,\\
 (T_N^{-1})_{rj}=\binom jr(-c)^{j-r}i^{-j}\quad(r\le j),
 \end{gathered}
\end{equation}
with zero entries below its triangular range, the original physical
unit in $J_{S,N}$ and the permutation $\Pi_N$ placing even powers
first. Expansion of $(c+iy)^j$ and $i^{-j}(S-c)^j$ proves both
coefficient maps. Keep $A=M_S$, the original class $x$ in the
packet, and the original observation kernel $K=\ker\Lambda$.
Set
\[
 W=(JJ^*)^{-1},\qquad
 \mathfrak A=\norm A_W\norm x_W^2,\qquad
 U=\sum_{b=0}^1\sum_{i=0}^{n_b}
 \frac{2A_kC_\Gamma[2(2i+b)]!}{\alpha^{2(2i+b)+1}}.
\]
The original upper envelope gives $H\preceq UI$, and its full
coefficient trace is bounded by
\begin{equation}\tag{MI19}\label{mi:U}
 U\le \overline U_N:=
 \frac{2A_kC_\Gamma(N+1)(2N)!}{\alpha}.
\end{equation}
There are $N+1$ diagonal terms and each has $2i+b\le N$;
the factors $\alpha^{-2(2i+b)}\le1$ prove this finite inequality.

Real source-moment midpoints give a Hermitian $\widehat H$ with
$\norm{H-\widehat H}\le d\varepsilon$. Put
$\theta=d\varepsilon/h_{\min}\le1/4$ and
$H_\pm=\widehat H\pm d\varepsilon I$.
Direct quadratic-form estimates give
\[
 H_-\succeq(1-2\theta)h_{\min}I,\quad
 H_-\preceq H\preceq H_+,\quad
 H_+\preceq(1-2\theta)^{-1}H_-.
\]
For every actual onto map $J$ and positive matrix $D$, completing
the square on the full fibre $Jb=u$ proves
\[
 \min_{Jb=u}b^*Db=u^*(JD^{-1}J^*)^{-1}u,\quad
 b_{\min}=D^{-1}J^*(JD^{-1}J^*)^{-1}u.
\]
The formula includes all of $\ker J$. Thus the original attained
metrics $G_0=(JH_-^{-1}J^*)^{-1}$ and
$G=(JH^{-1}J^*)^{-1}$ obey
$\frac12h_{\min}W\preceq G_0\preceq G
\preceq(1-2\theta)^{-1}G_0$ and $G_0\preceq UW$.

For clarity we reproduce the quantitative product transport rather
than replacing its phase by a norm comparison. Let $P_0,P$ be the
orthogonal projections onto the same $K$ in $G_0,G$, respectively.
Put $a=P_0x$, $b=(I-P_0)x$, $u=\norm a_{G_0}$, $v=\norm b_{G_0}$,
$e=(P-P_0)x$, and $M_A=\norm A_{G_0}$.
If $\kappa=(1-2\theta)^{-1}$, then
\[
 \norm e_{G_0}\le sv,\quad
 s=\frac{\kappa-1}{2\sqrt\kappa}
  =\frac{\theta}{\sqrt{1-2\theta}},\qquad e\in K.
\]
To prove this, write $G_0^{-1}G=\left[\begin{smallmatrix}B&X\\X^*&D
\end{smallmatrix}\right]$ in the $G_0$-orthogonal decomposition
$K\oplus K^{\perp_{G_0}}$. Orthogonality gives the off-diagonal
projection block $B^{-1}X$. The bound
$E^2\preceq(\kappa+1)E-\kappa I$ for $I\preceq E\preceq\kappa I$
gives
$B^{-1}XX^*B^{-1}\preceq
(\kappa+1)B^{-1}-\kappa B^{-2}-I
\preceq(\kappa-1)^2/(4\kappa)I$.
The last scalar maximum is attained at $2\kappa/(\kappa+1)$.

Define the unchanged complex quantities
$z=\langle Px,A(I-P)x\rangle_G$,
$w=\langle(I-P)x,APx\rangle_G$ and
$z_0=\langle a,Ab\rangle_{G_0}$,
$w_0=\langle b,Aa\rangle_{G_0}$, using conjugate-linearity first.
Expand $Px=a+e$, $(I-P)x=b-e$ and $G=G_0E$.
Both absolute differences are bounded by
\[
 M_A\{sv^2+suv+s^2v^2+
           (\kappa-1)(u+sv)\sqrt{1+s^2}\,v\}.
\]
This retains the three difference terms from expanding the two
vectors and the fourth from $E-I$. Since $\theta\le1/4$,
$s\le\sqrt2\theta$, $s^2\le2\theta^2$,
$\kappa-1\le4\theta$ and
$\sqrt{1+s^2}\le3/(2\sqrt2)$.
The braces are at most
$\theta[4\sqrt2uv+(2+\sqrt2)v^2]\le6\theta(u^2+v^2)$;
the last inequality is
$6(u-\sqrt2v/3)^2+(8/3-\sqrt2)v^2\ge0$.
Put $C=2\sqrt{U/h_{\min}}\,U\mathfrak A$.
The proved metric comparisons give
$M_A(u^2+v^2)\le C$ and
$|z_0|,|w_0|\le C/2$. Expanding the complete complex product,
including its quadratic error, yields
\begin{equation}\tag{MI20}\label{mi:current}
 |\overline zw-\overline{z_0}w_0|
 \le 6C^2\theta+36C^2\theta^2\le15C^2\theta,\qquad
 C^2=\frac{4U^3\mathfrak A^2}{h_{\min}}.
\end{equation}
Thus, for $\mathfrak A>0$ and requested absolute error $t_*>0$,
\begin{equation}\tag{MI21}\label{mi:epscurrent}
 \varepsilon\le
 \min\left\{\frac{h_{\min}}{4d},\
       \frac{t_*h_{\min}^2}{60dU^3\mathfrak A^2}\right\}.
\end{equation}
When $\mathfrak A=0$, positive definiteness of $W$ gives $A=0$
or $x=0$ and all four pairings vanish; only the first positivity
guard is needed. No class is recomputed from midpoint moments.

\section{Four original determinant signs and their cheaper budget}

Let the four original cutoffs be $N_i$ with their fixed original
restriction or attained quotient maps and resulting positive
metrics $Q_i$, of ranks $r_i$. Retain signs
$s_1=s_2=1$, $s_3=s_4=-1$.
At endpoint $i$ let $d_i$ and $h_i$ be the source block size and
improved source floor. For $r_i>0$ choose
\begin{equation}\tag{MI22}\label{mi:detbudget}
 \theta_i=\min\{1/4,t_*/(8r_i)\},\qquad
 \varepsilon_i=\theta_i h_i/d_i.
\end{equation}
Rank-zero determinants equal one. The same source argument and
same-fibre minimum above give
$Q_i^-\preceq Q_i\preceq Q_i^+
\preceq(1-2\theta_i)^{-1}Q_i^-$.
Restriction is an exact congruence; each subsequent onto quotient
uses the displayed attained-minimum formula on its complete fibre.
Consequently, with $e_i=\log\det Q_i-\log\det Q_i^-$,
\[
 0\le e_i\le-r_i\log(1-2\theta_i)
 \le4r_i\theta_i\le t_*/2.
\]
The inequality $-\log(1-2\theta)\le4\theta$ on $[0,1/4]$
follows by integrating $2/(1-2u)\le4$.
The positive pair bounds $e_1+e_2-e_3-e_4$ above and the negative
pair bounds it below. Therefore
\begin{equation}\tag{MI23}\label{mi:four}
 \left|\sum_{i=1}^4s_i\log\det Q_i
       -\sum_{i=1}^4s_i\log\det Q_i^-\right|\le t_*.
\end{equation}
There is no extra condition number for a fixed observation map
in this exact same-fibre comparison. Certified rounding of that map
and of subsequent arithmetic is a separate numerical error already
retained in \cite{oldNP}. To use one common moment sequence across
endpoints, choose $\varepsilon=\min_{i:r_i>0}\varepsilon_i$ and
$N_{\max}=\max_iN_i$ in \eqref{mi:tau}. If all ranks vanish the
return is exactly zero and no determinant tolerance is required.

\section{The evaluated scales of these explicit choices}

Take fixed-packet $h$ and a proportional window
$c_0q\le N\le c_1q$, with fixed positive $c_0,c_1$;
for four endpoints use the corresponding common bounds.
Equations \eqref{mi:log} and \eqref{mi:U}, together with
$\log((2N)!)\le2N\log(2N)$, give
\[
 \log^+(h_{\min}^{-1})=O_h(\sqrt q),\qquad
 \log^+U\le2N\log(N+2)+O_h(N+k).
\]
Choose the largest admitted error in \eqref{mi:epscurrent}, and
then the largest admitted $\tau$ in \eqref{mi:tau}. Taking
logarithms of those explicit minima gives
\begin{equation}\tag{MI24}\label{mi:currentcost}
 \log^+(\tau^{-1})
 \le8N\log(N+2)+2\log^+\mathfrak A+\log^+(t_*^{-1})
       +O_h(N+k+\sqrt N).
\end{equation}
This retains the actual finite-map factor $\mathfrak A$.
It assigns no polynomial bound to that factor.
For \eqref{mi:detbudget}, if $r_i\le c_2q$ with fixed $c_2$,
the largest common permissible moment error instead satisfies
$\log^+(\varepsilon^{-1})
=O_h(\sqrt q+\log(q+2)+\log^+(t_*^{-1}))$.
Its raw-moment tolerance obeys the stronger bound
\begin{equation}\tag{MI25}\label{mi:detcost}
 \log^+(\tau^{-1})
 \le2N_{\max}\log(N_{\max}+2)+\log^+(t_*^{-1})
       +O_h(N_{\max}+k+\sqrt q).
\end{equation}
These estimates describe the stated sufficient choices, not
unnecessarily smaller tolerances. Exact determinant centres,
the original correlated current values and the original complete
action remain the same mathematical targets.

\begin{thebibliography}{9}
\bibitem{Dunster} T.~M.~Dunster, \emph{Legendre and Related Functions},
NIST Digital Library of Mathematical Functions, Chapter 14,
Rodrigues formula 14.7.13, version 1.2.8 (15 September 2026).
\url{https://dlmf.nist.gov/14.7.E13}. Original equation TeX read.
The chapter credits earlier work including I.~A.~Stegun and
F.~W.~J.~Olver; their books were not separately read for this note.
\bibitem{DLMF18} T.~H.~Koornwinder, R.~Wong, R.~Koekoek,
R.~F.~Swarttouw and W.~P.~Reinhardt,
\emph{Orthogonal Polynomials}, NIST Digital Library of Mathematical
Functions, Chapter 18, 18.12.11, version 1.2.8.
\url{https://dlmf.nist.gov/18.12.E11}. Original equation TeX read.
\bibitem{BT} Split-Zero programme,
\emph{Actual boundary tilts and uniform convolution tails},
complete provider \texttt{BOUNDARY\_TILT\_COMPLETE.tex}, BT1--23.
The exact BT12 constant and BT14--16 estimates are used in MI15.
\bibitem{oldNP} Split-Zero programme,
\emph{Explicit native moment precision, inverse margins, and the
original signed receiver}, 20 September 2026,
\texttt{NATIVE\_MOMENT\_PRECISION.tex}, original envelopes,
weighted systems, quadratic-expression guards and finite moments.
\bibitem{MC} Split-Zero programme,
\emph{An explicit original-native moment tolerance for the complex
mixed product}, 20 September 2026,
\texttt{NATIVE\_MOMENT\_PHASE\_COMPOSITION.tex}, MC1--16.
MI18--21 retain and reprove its full constant-$15$ product estimate
with the improved native floor.
\bibitem{NR} Split-Zero programme,
\emph{A finite auxiliary-degree bound for the original native
resolvent}, 20 September 2026,
\texttt{NATIVE\_RESOLVENT\_DEGREE\_RATE.tex}, NR19 and NR26.
MI13 supplies a second certified inverse-evaluation bound to the
same original degree prescription.
\end{thebibliography}
\end{document}
